\section{Contexto macro e o Conditional Autoencoder}
\label{sec:cond-macro}

Até aqui o modelo olha só para o \emph{ativo} (preço e volume). Mas uma queda pode ter duas causas
muito diferentes: um problema \textbf{do próprio ativo} (uma fraude, um \emph{rombo} de liquidez)
ou um \textbf{choque do mercado inteiro} (uma crise global). Para distinguir, damos à rede um
\emph{contexto macroeconômico}: dólar (USD/BRL), VIX (o ``índice do medo''), Selic e IPCA.

\subsection{Por que não basta jogar a macro junto}
\label{sec:cond-juntar}

\begin{intuicao}
Se você empilha a macro no mesmo bloco que a rede tenta \emph{reconstruir}, acontece uma injustiça:
o IPCA quase não muda dentro de uma janela de 30 dias (é mensal), então reconstruí-lo é trivial ---
a rede ``acerta'' sem esforço. O erro de reconstrução passa a vir só de preço/volume, e a macro
vira decoração: está lá, mas não influencia nada.
\end{intuicao}

A causa é a perda \textbf{MSE}: ela é dominada pelas variáveis que mais variam. Uma série quase
constante contribui com erro $\approx 0$ e é \emph{ignorada} pelo treino.

\subsection{A ideia: macro como contexto, não como alvo}
\label{sec:cond-ideia}

\begin{definicao}[Conditional Autoencoder]
Um autoencoder em que parte da entrada serve apenas para \emph{condicionar} (dar contexto) a
codificação, mas \textbf{não} faz parte do que se tenta reconstruir. Aqui: o \emph{encoder} vê
preço/volume \emph{e} macro; o \emph{decoder} reconstrói \textbf{só} preço/volume; a perda mede só
o erro de preço/volume.
\end{definicao}

\begin{intuicao}
A macro fica como ``pano de fundo'': ela molda a representação interna (o gargalo), mas a rede nunca
é cobrada por reconstruí-la. Assim a macro \emph{informa} sem ser ignorada nem poluir a conta do
erro.
\end{intuicao}

\subsection{Separando idiossincrático de sistêmico}
\label{sec:cond-regime}

Com dois números por janela conseguimos a distinção:
\begin{itemize}[noitemsep]
  \item \textbf{erro de Preço/Volume} --- o quanto o \emph{ativo} está estranho (vem da
        reconstrução);
  \item \textbf{estresse Macro} --- o quanto o \emph{mercado} se mexeu (medido direto da macro, pela
        amplitude dela na janela --- não por reconstrução).
\end{itemize}

\begin{exemplo}
Ativo muito estranho $+$ macro parada $\Rightarrow$ \textbf{idiossincrático} (o problema é dele:
a fraude da Americanas). \\
Ativo muito estranho $+$ macro também agitada $\Rightarrow$ \textbf{sistêmico} (a crise é de todos:
o \emph{crash} da COVID, com o VIX nas alturas).
\end{exemplo}

\subsection{Cuidado com o tempo: data de publicação, não de referência}
\label{sec:cond-leak}

\begin{intuicao}
O IPCA de \emph{março} só é \emph{divulgado} em meados de \emph{abril}. Se você ``souber'' o IPCA de
março já em 1º de março, está espiando o futuro (Seção~\ref{sec:vazamento}). Por isso cada indicador
é datado pela sua \textbf{data de publicação}, e o preenchimento é só ``para frente'' (\emph{ffill}):
o dado existe para a rede a partir do dia em que virou público --- nunca antes.
\end{intuicao}

\begin{naprat}
Validado no notebook \texttt{12}: o \emph{crash} da COVID/2020 foi classificado \textbf{sistêmico}
e a fraude da Americanas/2023, \textbf{idiossincrático} --- a separação que o método buscava. O
contexto macro vem sobretudo das séries \emph{diárias} (USD/BRL, VIX); as mensais (Selic, IPCA) são
quase constantes na janela e contribuem pouco. Decisão registrada no ADR-0012.
\end{naprat}
